% !TEX program = lualatex
\documentclass[a4paper,10pt,twocolumn]{article}
\usepackage{kaitoushu}
\renewcommand{\baselinestretch}{1.2}
\lhead{応用数学 問題集}
\problemrange{問90〜92}

\begin{document}
\thispagestyle{fancy}

\begin{center}
  {\Large \textbf{応用数学　2章　第4回　Basic　解答}}
\end{center}
\vspace{0.5em}

\noindent\textbf{\large【Basic】}
\vspace{0.3em}

\mondai{90}{
次の関数の逆ラプラス変換を求めよ．
(1) $\dfrac{1}{s^2+4s+4}$ \quad
(2) $\dfrac{1}{s^2-4s+3}$ \quad
(3) $\dfrac{1}{s^2-6s+10}$
\hfill {\scriptsize [教p.59(問17)]}
}

\kotae{
(1) $\dfrac{1}{(s+2)^2}$ より $\mathcal{L}^{-1} = te^{-2t}$

(2) $\dfrac{1}{(s-1)(s-3)} = \dfrac{1}{2}\!\left(\dfrac{1}{s-3}-\dfrac{1}{s-1}\right)$ より $\mathcal{L}^{-1} = \dfrac{1}{2}(e^{3t}-e^t)$

(3) $\dfrac{1}{(s-3)^2+1}$ より $\mathcal{L}^{-1} = e^{3t}\sin t$
}

\mondai{91}{
次の関数の逆ラプラス変換を求めよ．
(1) $\dfrac{s+1}{s^2+6s+9}$ \quad
(2) $\dfrac{s}{s^2-4s+13}$
\hfill {\scriptsize [教p.59(問18)]}
}

\kotae{
(1) $\dfrac{s+1}{(s+3)^2} = \dfrac{(s+3)-2}{(s+3)^2} = \dfrac{1}{s+3} - \dfrac{2}{(s+3)^2}$

$\mathcal{L}^{-1} = e^{-3t} - 2te^{-3t}$

(2) $\dfrac{s}{(s-2)^2+9} = \dfrac{(s-2)+2}{(s-2)^2+9}$

$\mathcal{L}^{-1} = e^{2t}\cos 3t + \dfrac{2}{3}e^{2t}\sin 3t$
}

\mondai{92}{
次の関数の逆ラプラス変換を求めよ．
(1) $\dfrac{2s^2-5s-6}{(s+2)(s-1)(s-2)}$ \quad
(2) $\dfrac{1}{s^2(s+1)}$
\hfill {\scriptsize [教p.60(問19)]}
}

\kotae{
(1) 部分分数分解：$\dfrac{2s^2-5s-6}{(s+2)(s-1)(s-2)} = \dfrac{A}{s+2}+\dfrac{B}{s-1}+\dfrac{C}{s-2}$

$A=\dfrac{8+10-6}{(-3)(-4)}=1$，$B=\dfrac{2-5-6}{(3)(-1)}=3$，$C=\dfrac{8-10-6}{(4)(1)}=-2$

$\mathcal{L}^{-1} = e^{-2t}+3e^t-2e^{2t}$

(2) $\dfrac{1}{s^2(s+1)} = \dfrac{A}{s}+\dfrac{B}{s^2}+\dfrac{C}{s+1}$

$A=-1$，$B=1$，$C=1$ より $\mathcal{L}^{-1} = -1+t+e^{-t}$
}

\end{document}